\begin{center}\begin{tikzpicture}[font=\small,>=stealth,line join=round]
\draw[gray!40,rounded corners=3pt] (0,0) rectangle (14.2,8.3);\node at (7.1,7.700000000000001) {例5：立方与正弦都把相反数保留下来};
\begin{scope}[shift={(4.1,3.9499999999999997)},x=1.5cm,y=1.6428571428571428cm]\draw[->,gray] (-2,0)--(2,0) node[right]{$x$};\draw[->,gray] (0,-1.4)--(0,1.4) node[above]{$y$};\draw[blue,thick,solid,domain=-1.8:1.8,samples=180] plot(\x,{sin(deg((\x)^3))});\end{scope}\node[anchor=west,align=left,text width=5.0cm] at (8.4,4.0) {$F=\sin(x^3)$\\[10pt]内层：$(-x)^3=-x^3$\\[10pt]外层：$\sin(-u)=-\sin u$\\[10pt]整体为奇函数};
\node[text width=12.6cm,align=center] at (7.1,.65) {起伏由真实公式决定；奇偶性关心原点对称，不要求单调。};
\end{tikzpicture}\end{center}
